\clearpage
\section{Sharing completed query results}
\label{sec:bioseq-result-cache}

An alignment query may run for hours while returning a small result. If two
sessions submit the same request, running the alignment twice wastes work.
The separate \texttt{bioseq\_cache} extension coordinates these requests and
stores completed results. This differs from the two caches used when reading
sequences: a lazy read session can reuse fetched fragments, and the column
cache reuses decompressed storage chunks within a backend. The result cache
can serve another execution without repeating alignment.

The extension provides two forms of reuse:
\begin{itemize}
  \item \textbf{Exact results.} An identical eligible request can wait for a
    running producer or read its completed result, including repeated rows
    and their stored order.
  \item \textbf{Reusable fragments.} A provider can describe which input
    rows a completed computation covers. A later request can read suitable
    results and compute only the missing part.
\end{itemize}
The generic cache implements both forms. The current BioSeqDB adapter
provides exact result reuse; it does not yet provide the coverage and row
identity callbacks needed for reusable fragments. The range example below
explains the generic mechanism and the contract such an adapter must meet.

\subsection{One producer, waiting readers}

An entry is one attempt to compute a request over compatible data versions.
The \emph{producer} owns that attempt. Other executions either wait for it
or read a completed entry. To \emph{seal} the result is to finish writing it
and close it to further rows. Figure~\ref{fig:cache-lifecycle} shows the
successful path and the two ways in which an attempt stops being reusable.

\begin{figure}[htbp]
\centering
\begin{tikzpicture}[x=1cm,y=1cm,font=\footnotesize]
  \draw[scan arrow,black!55] (0,3.0) -- (11.7,3.0)
    node[right]{time};
  \node[anchor=east] at (0,2.25) {producer};
  \node[anchor=east] at (0,.85) {waiter};
  \node[scan input,minimum width=1.2cm] (claim) at (.9,2.25) {claim};
  \node[scan input,minimum width=3.6cm] (build) at (3.7,2.25)
    {compute and store rows};
  \node[scan output,minimum width=1.8cm] (seal) at (6.9,2.25)
    {seal and check};
  \node[scan output,minimum width=1.6cm] (complete) at (9.2,2.25)
    {complete};
  \draw[scan arrow] (claim) -- (build);
  \draw[scan arrow] (build) -- (seal);
  \draw[scan arrow] (seal) -- (complete);
  \node[scan box,minimum width=1.25cm] (lookup) at (2.25,.85) {lookup};
  \draw[thick,dashed,scanBlue] (lookup.east) -- (7.65,.85)
    node[midway,above=3pt]{wait; no partial rows};
  \node[scan output,minimum width=1.5cm] (recheck) at (9.2,.85)
    {check again};
  \node[scan output,minimum width=1.25cm] (read) at (11.2,.85) {read};
  \draw[scan arrow] (7.65,.85) -- (recheck);
  \draw[scan arrow,scanGreen] (complete) -- (recheck)
    node[midway,right=2pt]{wake};
  \draw[scan arrow] (recheck) -- (read);
  \node[scan hit,minimum width=2.6cm] (failed) at (3.7,-.65)
    {failed attempt};
  \draw[thick,scanOrange] (build.south) -- (3.7,1.45);
  \draw[thick,scanOrange] (seal.south) -- (6.9,1.45) -- (1.2,1.45);
  \fill[scanOrange] (3.7,1.45) circle[radius=1.1pt];
  \draw[scan arrow,scanOrange] (1.2,1.45) |- (failed.west);
  \node[anchor=east,text=scanOrange,align=right] at (1.05,1.45)
    {error or\\early stop};
  \node[scan hit,minimum width=3.25cm] (invalid) at (9.7,-.65)
    {expired or invalidated};
  \draw[scan arrow,scanOrange] (complete.east) -- (12.15,2.25)
    |- (invalid.east);
  \node[scan box,minimum height=.5cm] (removed) at (6.75,-1.55)
    {removed when unused};
  \draw[scan arrow,black!60] (failed.south) |- (removed.west);
  \draw[scan arrow,black!60] (invalid.south) |- (removed.east);
\end{tikzpicture}
\caption{Two executions share one computation. The producer can return
its own rows while building the entry. A waiter reads only after completion
and validation. Failed, expired, or invalidated entries are removed after
their users release them.}
\label{fig:cache-lifecycle}
\end{figure}

The lifecycle has four main rules:
\begin{itemize}
  \item \textbf{Claim briefly, compute without the metadata lock.}
    A short shared-state update chooses one producer for a matching request
    and data version. The producer does not hold this lock throughout the
    alignment. Waiters sleep and periodically check cancellation, producer
    liveness, and data validity.
  \item \textbf{Publish only complete work.} Reaching the child's end of
    output allows the producer to seal the stored result, check dependencies,
    and mark the entry complete. An early cursor close cannot publish the
    rows collected so far as a complete result.
  \item \textbf{Validate before reading.} A wake-up means that the entry
    changed. The reader checks it again and keeps a reference that prevents
    removal while it reads. The retention time starts at completion.
  \item \textbf{Release on every exit.} Canceling a waiter releases that
    waiter's reference and leaves the producer running. A failed or canceled
    producer wakes its waiters, which can run the retained ordinary plan.
    A later request can start a new attempt. Waiting is declined when
    executions could end up waiting for one another.
\end{itemize}

Matching SQL text alone is insufficient. Exact identity includes the
analyzed query, parameter values, result types, relevant execution settings,
and the database and security context. Data validity is checked separately
against the query's snapshot and recorded dependency versions. For example,
an identical query after a committed change to its dependencies must not
return an invalidated result.
Permissions still undergo normal PostgreSQL checks. Unsupported or invalid
cache entries leave the ordinary query plan available.
Section~\ref{sec:pg-storage-request-coordination} describes result storage
and its memory and file limits.

\subsection{Describe coverage, then compute the difference}

Suppose the expensive operation and its settings remain fixed, while the
user changes a metadata filter. A reusable entry stores both its results
and the predicate describing the input rows it processed. A \emph{cached
source} is a completed entry selected for reuse. This is the idea
of \emph{semantic caching}: describe the cached region, read the useful
part, and compute a remainder.\footnote{Dar, Franklin, J\'onsson,
Srivastava, and Tan, \href{https://www.vldb.org/conf/1996/P330.PDF}
{\emph{Semantic Data Caching and Replacement}}, VLDB 1996, pp.~330--341.}
Coverage records completed work even when some input rows produced no hits.
The absence of a saved hit alone cannot prove that its input was searched.

For example, select input rows with $40\leq id<100$. Completed sources cover
$0\leq id<60$ and $50\leq id<80$. We read $[40,60)$ from the first source,
$[60,80)$ from the second, and compute $[80,100)$. These ranges select
metadata row identifiers, not positions within a biological sequence.

We can state the rule in tuple relational calculus. Fix one compatible
database snapshot, search query, alignment settings, and requested match
fields that affect which hits are emitted. Define:
\begin{itemize}
  \item $T(t)$: $t$ is an input row whose identity is retained;
  \item $A(t,h)$: the operation emits hit occurrence $h$ for row $t$;
    $h$ includes a distinct identity even when two hits have equal displayed
    values;
  \item $p(t)$: the new SQL predicate $P$ evaluates to \texttt{TRUE} on $t$;
  \item $c_i(t)$: cached source $i$ passed validation and its recorded
    predicate $P_i$ evaluates to \texttt{TRUE} on $t$.
\end{itemize}
Order the selected sources from $1$ to $m$. Source $i$ supplies branch $B_i$;
the relation $D$ contains the input rows that still need computation:
\begin{align}
 B_i &= \{(t,h)\mid T(t)\land A(t,h)\land p(t)\land c_i(t)
                     \land \forall j<i:\neg c_j(t)\},
                     \label{eq:cache-branch}\\
 D &= \{t\mid T(t)\land p(t)\land \forall i\in\{1,\ldots,m\}:\neg c_i(t)\}.
                     \label{eq:cache-difference}
\end{align}
Each requested hit is returned from the first valid source covering its
input, or computed from $D$. Thus
\begin{equation}
 \{(t,h)\mid T(t)\land A(t,h)\land p(t)\}
   = B_1\uplus\cdots\uplus B_m
     \uplus\{(t,h)\mid t\in D\land A(t,h)\}.
 \label{eq:cache-result}
\end{equation}
Here $\uplus$ joins disjoint sets of occurrences. The operation is evaluated
only on missing inputs $D$; $A$ is a description of the result, not an
instruction to compute every possible hit first. Removing the hidden
identities afterwards preserves repeated SQL rows. SQL \texttt{UNION ALL}
has the required concatenation behavior; \texttt{UNION} could remove
legitimate repeated rows.\footnote{PostgreSQL 17,
\href{https://www.postgresql.org/docs/17/queries-union.html}
{Combining Queries}.}

\subsection{An overlapping range example}

Figure~\ref{fig:cache-ranges} shows the same example, including the overlap
between the two sources.

\begin{figure}[htbp]
\centering
\begin{tikzpicture}[x=.10cm,y=1cm,font=\footnotesize]
  \foreach \x in {0,40,50,60,80,100} {
    \draw[black!18] (\x,-.15) -- (\x,3.1);
    \node[anchor=north] at (\x,-.2) {$\x$};
  }
  \node[anchor=east] at (-4,2.75) {request};
  \draw[fill=black!8,draw=black!70] (40,2.52) rectangle (100,2.98);
  \node[anchor=east] at (-4,1.95) {source 1};
  \draw[fill=scanBlue!15,draw=scanBlue] (0,1.72) rectangle (60,2.18);
  \node[anchor=east] at (-4,1.15) {source 2};
  \draw[fill=scanGreen!16,draw=scanGreen] (50,.92) rectangle (80,1.38);
  \fill[black!20] (50,.92) rectangle (60,1.38);
  \node[anchor=east] at (-4,.35) {return from};
  \draw[fill=scanBlue!20,draw=scanBlue] (40,.12) rectangle (60,.58);
  \draw[fill=scanGreen!20,draw=scanGreen] (60,.12) rectangle (80,.58);
  \draw[fill=scanOrange!20,draw=scanOrange] (80,.12) rectangle (100,.58);
  \node at (50,.35) {source 1};
  \node at (70,.35) {source 2};
  \node at (90,.35) {compute};
  \draw[scan arrow,black!55] (0,-.65) -- (103,-.65)
    node[right]{row $id$};
\end{tikzpicture}
\caption{The new request uses source 1 for $[40,60)$, source 2 for
$[60,80)$, and computes $[80,100)$. The gray overlap in source 2 is omitted
because source 1 already supplies those occurrences. Cached rows outside
the new request are also omitted.}
\label{fig:cache-ranges}
\end{figure}

With both sources valid, the missing inputs have this equivalent SQL
description:
\begin{lstlisting}[style=bioseq-search-sql,
  caption={Input rows on which the expensive operation is still needed.},
  label={lst:cache-missing-inputs}]
SELECT id, seq
FROM sequences
WHERE id >= 40 AND id < 100
  AND (id >=  0 AND id < 60) IS NOT TRUE
  AND (id >= 50 AND id < 80) IS NOT TRUE;
\end{lstlisting}
The cache builds PostgreSQL expression nodes for this condition; it does
not construct and parse a new SQL string. A provider must place the
condition before the expensive operation. The listing describes the input
selection and does not introduce a new public search function.

Three details make this transformation correct:
\begin{itemize}
  \item \textbf{Handle SQL NULLs.} A \texttt{WHERE} clause accepts only
    \texttt{TRUE}. For a nullable predicate, \texttt{NOT P} also becomes
    unknown when $P$ is unknown. \texttt{P IS NOT TRUE} includes both false
    and unknown, so it keeps every input excluded from the earlier
    cache.\footnote{PostgreSQL 17,
    \href{https://www.postgresql.org/docs/17/functions-comparison.html}
    {Comparison Functions and Operators}, Table~9.2.}
  \item \textbf{Restore invalid coverage.} All source checks finish before
    output begins. If source 1 fails validation, $c_1$ is false: source 2
    supplies $[50,80)$ and $D$ becomes $[40,50)\cup[80,100)$. An internal
    Boolean parameter enables each exclusion only for an accepted source.
  \item \textbf{Require enough information.} A reusable provider must keep
    input and hit identities, prove that results have compatible meaning,
    and supply columns needed to test the new predicate. Missing metadata
    may be fetched from the original rows under the validated snapshot.
    A result containing coordinates alone cannot supply an uncomputed CIGAR.
\end{itemize}

Coverage is taken at a complete expensive-operation boundary. Joins,
aggregation, and ordering above that boundary run again on the reconstructed
rows. If an upper \texttt{LIMIT} stops the boundary early, its partial rows
cannot become a complete coverage entry, even when the final limited query
has a complete exact result. These rules allow reuse to save computation
while keeping PostgreSQL responsible for the surrounding query.
\label{sec:bioseq-result-cache-end}
